\chapter{破坏经济秩序犯罪}

% Source PDF physical page 333; see audit/source_coverage.json.
\begin{EditionOpeningPanel}
复习策略：由于知识点比较散乱，所以尽量总结规律，通过对比、图表来记忆。
\end{EditionOpeningPanel}

\section{生产、销售伪劣商品罪}

[问题】电影《我不是药神》中，男主进口印度仿制药，有疗效，但没有我国药监部门的批号。放到现今，该如何处理？（药神案）

\subsection*{一、生产、销售伪劣产品罪}

第140条生产者、销售者在产品中掺杂、掺假，以假充真，以次充好或者以不合格产品冒充合格产品，\CJKunderwave{销售金额五万元以上}的，构成本罪。

1.本罪的“销售”应作缩小解释，要求销售行为具有“推向市场”的性质，因为本罪是破坏市场秩序的犯罪，保护的法益是广大购买者（不特定人）的合法权益。

例1，甲欠乙10万元债务，甲将价值1万元的低档白酒冒充价值10万元的高档白酒，交给乙，用来抵偿债务。甲的行为不具有“推向市场”的性质，没有破坏市场秩序，不构成销售伪劣产品罪，但构成诈骗罪。（2024年试题）

例2，乙公司向甲工厂购买办公用品，双方合同约定了产品质量标准、若未达到该质量标准则应承担的违约责任。合同履行中，甲工厂提供了质量不合格产品。乙公司索赔，甲工厂承担了违约赔偿责任。甲工厂不构成销售伪劣产品罪。该案按照民事纠纷处理即可。若给甲工厂定本罪，则会导致大量民事纠纷变成刑事犯罪，会导致大量企业不敢做生意了。

2.销售金额5万元

（1）达到5万元，构成既遂。

（2）未达到5万元的，是未遂，但轻微的未遂不值得处罚。司法解释规定：（销售金额\ensuremath{\times}3）+货值金额\ensuremath{\ge}15万元，则以生产、销售伪劣产品罪（未遂）定罪处罚。

% Source PDF physical page 334; see audit/source_coverage.json.
例1，销售金额是0，货值金额（库存额）是15万元，以生产、销售伪劣产品罪（未遂）定罪处罚。（2010年第15题）

例2，销售金额是2万元，货值金额是9万元。（2\ensuremath{\times}3）+9=15万元，以生产、销售伪劣产品罪（未遂）定罪处罚。

3.本罪与本节其余九个具体罪名的关系

（1）\CJKunderwave{第149条第1款}：生产、销售九个具体罪名的产品，不构成该罪的，如果销售金额在5万元以上，定生产、销售伪劣产品罪。

（2）\CJKunderwave{第149条第2款}：如果行为同时构成九个具体罪名和生产、销售伪劣产品罪，择一重罪论处。

4.本节罪名与诈骗罪：销售伪劣产品，同时触犯诈骗罪的，想象竞合，择一重罪论处。

5.本节罪名均可以由单位构成。

\subsection*{二、食品犯罪}

\InsertDiagram{figures/manual_table_p334_1.tex}

[注意］生产、销售有毒、有害食品罪与投放危险物质罪可以想象竞合。例如，甲用工业酒精冒充白酒出售，导致数人中毒，同时触犯两罪，想象竞合，择一重罪论处。（2022年试题）

% Source PDF physical page 335; see audit/source_coverage.json.
\subsection*{三、药品犯罪}

\InsertDiagram{figures/manual_table_p335_1.tex}

[练习]某药厂明知药品已经过期，仍然作为捐献的药品捐出，造成严重后果。药厂不构成提供劣药罪，因为药厂是药品的生产单位，不是使用单位。\footnote{少数说认为，药厂构成提供劣药罪，因为将药厂解释为药品使用单位，属于扩大解释。}（2022年试题）

% Source PDF physical page 336; see audit/source_coverage.json.
\InsertDiagram{figures/manual_p336_1.tex}

[练习]甲没有获得药品批准文件，销售国外合法的仿制药。该仿制药对治疗丙肝有疗效。甲构成：销售假药罪？妨害药品管理罪？非法经营罪？答案：甲不构成犯罪。

\section{走私犯罪}

[问题]甲误将淫秽光盘当作普通光盘走私入境。对甲如何处理？（走私淫秽光盘案）

“\CJKunderwave{走私}”有两种含义：

1.\CJKunderwave{针对国家允许进出口的物品，偷逃海关关税而进出口}。对此的罪名是走私普通货物、物品罪。该罪实际上是海关领域的逃税罪。由于偷逃的关税主要是进口关税，所以该罪的对象主要是进口的货物、物品。

2.\CJKunderwave{针对国家禁止进出口的物品，违反海关制度实施进出口}。对此，有个兜底罪名“走私国家禁止进出口的货物、物品罪”，还有八个具体的罪名，如走私武器、弹药罪，走私假币罪等。

\InsertDiagram{figures/manual_p336_2.tex}

\subsection*{一、走私普通货物、物品罪}

1.变相走私（第154条）。这是指变相偷逃海关关税，主要方式是，将保税货物在境内销售牟利。

例如，某个“三来一补”企业，正常流程是，从国外进口原材料时，海关不收进口关税，但要求其将原材料加工成成品后必须出口到国外，出口时海关也不再给该企业出口退税。进口的这种原材料就是保税货物。该企业将该原材料在境内销售牟利，属于变相偷逃海关关税。

% Source PDF physical page 337; see audit/source_coverage.json.
[注意］一般的走私不要求具有非法牟利目的，而该种走私（出口转内销），法条加了“牟利”二字，要求具有非法牟利目的。例如，来料加工的原材料因为质量问题不可能复运出境或没有任何经济价值和必要性复运出境，然后在境内销售，未补缴关税。这种行为不构成走私罪。但是，由于走私普通货物、物品罪实际是海关领域的逃税罪，所以，未补缴的关税的数额达到逃税罪的标准的，以逃税罪论处。但是，这种逃税罪不能享受逃税罪的排除处罚的待遇，也即第201条第4款规定的“经税务机关下达追缴通知后，补缴应纳税款，·….不予追究刑事责任”，因为海关毕竞不是税务机关。（2017年第13题）

2.间接走私（第155条）。两种情形：第一，直接向走私人非法收购走私物品；第二，在境内运输、收购、贩卖走私物品。\CJKunderwave{注意}：对间接走私不再定掩饰、隐瞒犯罪所得罪。

\subsection*{二、走私禁止进出口的货物、物品}

本部分共九个罪名（第151条、第152条）。兜底罪名是走私国家禁止进出口的货物、物品罪。

\subsubsection*{（一）走私方向}

这九个罪，六个罪是双向禁止，禁止进出口。三个罪是单向禁止：

1.走私文物罪与走私贵重金属罪：只禁止出口，不禁止进口，这两个罪的走私方向是出口。

2.走私废物罪：只禁止进口，不禁止出口，该罪的走私方向是进口。

\subsubsection*{（二）包容评价}

\InsertDiagram{figures/manual_p337_1.tex}

\subsubsection*{（三）走私武器、弹药罪}

\InsertDiagram{figures/manual_table_p337_2.tex}

% Source PDF physical page 338; see audit/source_coverage.json.
\section{妨害对公司、企业的管理秩序罪}

\subsection*{一、非国家工作人员受贿罪（第163条）}

\CJKunderwave{这是指，公司、企业或者其他单位的工作人员}利用职务上的便利，索取他人财物或者非法收受他人财物，为他人谋取利益，数额较大的行为。人罪数额是3万元，与受贿罪相同。

\InsertDiagram{figures/manual_table_p338_1.tex}

\subsection*{二、对非国家工作人员行贿罪（第164条）}

1.对国家工作人员行贿，构成行贿罪（第389条）。两个行贿罪都要求为谋取不正当利益。两个行贿罪的人罪数额都是3万元。

2.消极的构成要件要素。行贿罪（第389条）中规定了消极的构成要件要素：“因被勒索给予国家工作人员以财物，没有获得不正当利益的，不是行贿”，而对非国家工作人员行贿罪中没有这种规定，但应当认为，因被勒索给予非国家工作人员以财物，没有获得不正当利益的，也不构成该罪。这是一种有利于被告人的类推解释。

3.对外国公职人员、国际公共组织官员行贿，构成对外国公职人员、国际公共组织官员行贿罪（第164条）。该罪要求谋取的是\CJKunderwave{不正当的商业利益}。

[提示】外国公职人员、国际公共组织官员受贿的，如果发生在我国境内，则构成我国的非国家工作人员受贿罪。例如，国际公共组织官员甲在中国境内收受乙的财物。甲构成非国家工作人员受贿罪，乙既触犯对非国家工作人员行贿罪，也触犯对国际公共组织官员行贿罪，二者是法条竞合关系，以后者论处。

\subsection*{三、为亲友非法牟利罪}

第166条国有公司、企业、事业单位的工作人员，利用职务便利，有下列情形之一致使国家利益遭受重大损失的，处三年以下有期徒刑或者拘役，并处或者单处罚金；··

\subsubsection*{（一）\CJKunderwave{将本单位的盈利业务交由自己的亲友进行经营}的；}

\subsubsection*{（二）以\CJKunderwave{明显高于市场的价格}从自己的亲友经营管理的单位\CJKunderwave{采购商品、接受服务}或者}

以\CJKunderwave{明显低于市场的价格}向自己的亲友经营管理的单位\CJKunderwave{销售商品、提供服务}的；

\subsubsection*{（三）从自己的亲友经营管理的单位\CJKunderwave{采购、接受不合格商品、服务}的。}

\CJKunderwave{其他公司、企业的工作人员违反法律、行政法规规定，实施前款行为，致使公司、企业}\CJKunderwave{利益遭受重大损失的，依照前款的规定处罚}。[第2款是《刑法修正案（十二）》增设的］

[提示]本罪与贪污罪可以形成想象竞合，择一重罪论处。例如，A国有公司原本直

% Source PDF physical page 339; see audit/source_coverage.json.
接从B公司购买电子产品，A公司的总经理王某让表弟赵某成立C公司，由C公司从B公司购买该产品，然后A公司以更高的价格向C公司购买，从而使C公司获利、A公司受损。也即，原来的交易流程是：B\ensuremath{\rightarrow}A，现在变成：B\ensuremath{\rightarrow}C\ensuremath{\rightarrow}A。这属于本罪的第二项行为类型。王某构成本罪。另外，A公司平白无故遭受了公共财产的损失，该损失转化为C公司的利润，等于C公司占有了该公共财产。就此，王某构成贪污罪。虽然王某自已没有占有该公共财产或利润，但是贪污罪所要求的非法占有目的，既包括为自己非法占有，也包括为第三人非法占有。王某同时构成本罪和贪污罪，想象竞合，择一重罪论处。

\subsection*{四、非法经营同类营业罪}

第165条国有公司、企业的董事、监事、高级管理人员，利用职务便利，自己经营或者为他人经营与其所任职公司、企业同类的营业，获取非法利益，数额巨大的，处三年以下有期徒刑或者拘役，并处或者单处罚金；数额特别巨大的，处三年以上七年以下有期徒刑，并处罚金。

其他公司、企业的董事、监事、高级管理人员违反法律、行政法规规定，实施前款行为，致使公司、企业利益遭受重大损失的，依照前款的规定处罚。[第2款是《刑法修正案（十二）》增设的］

\section{破坏金融管理秩序罪}

[问题］狗蛋隐瞒真相，使用假人民币向小芳兑换真美元。对狗蛋如何处理？（兑换美元案）

\subsection*{一、货币犯罪}

\subsubsection*{（一）伪造货币罪与变造货币罪（第170条、第173条）}

\InsertDiagram{figures/manual_table_p339_1.tex}

% Source PDF physical page 340; see audit/source_coverage.json.
\subsubsection*{（二） 出售、购买、运输假币罪（第 171 条)}

1.出售，是指有偿转让，条件是对方知晓是假币，否则构成使用假币罪（存在诈骗）。出售的方式既包括用假币换取真币，也包括用假币换取实物。

2.购买，是指有偿取得假币。\CJKunderwave{注意不能犯的情形}。例如，甲将冥币或白纸冒充假币出卖给乙。甲不定出售假币罪，而定诈骗罪。乙不成立购买假币罪，属于不能犯，不能定购买假币罪未遂。

3.运输，是指转移假币的存放地点。第一，这里的“运输”要求为了出售而运输。为了自己使用，在外地购买假币后，携带假币乘坐交通工具返回居住地的，不构成运输假币罪，而构成持有假币罪。第二，运输限于国内运输。出入国（边）境运输，构成走私假币罪。

\subsubsection*{（三）持有、使用假币罪（第172条）}

\CJKunderwave{使用}是指将假币作为真币投入流通领域。

1.\CJKunderwave{对人使用，要求对方不知情}。

（1）用假币进行交易，属于使用假币。例如，购物、消费、还债、赌博、兑换外币。（2015年第15题）

（2）将假币单向交付给对方，属于使用假币。例如，缴纳税款或罚款、冒充赃款上交、存入银行作为注册资本、赠送给不知情者。（2010年第13题）

[注意］没有交付使用而只是显示，不是使用假币。例如，将假币作为经济实力加以显示，不属于使用假币。

2.\CJKunderwave{对机器使用假币}。这里的机器是指自动售货机、自动存取款机。例如，向自动售货机投入假币换取商品，通过自动存取款机将假币存入银行，都属于使用假币。

3.使用假币罪与出售假币罪的区分标准：使用假币罪中，交易对方不知是假币；出售假币罪中，交易对方知道是假币。

例1，本节开头“兑换美元案”中，狗蛋同时触犯使用假币罪和诈骗罪，想象竞合，择一重罪论处。

例2，狗蛋使用假人民币向小芳兑换假美元，双方都知道都是假币。狗蛋属于出售假人民币，购买假美元，构成出售、购买假币罪。小芳属于出售假美元，购买假人民币，构成出售、购买假币罪。

例3，甲欺骗乙：“我将10万元假币卖给你，你给我1万元。”乙答应照办。乙收到后发现全部是白纸。甲对乙不构成出售假币罪，而构成诈骗罪既遂。乙不构成购买假币罪，因为甲欺骗乙，乙不可能买到假币。乙是甲的诈骗罪的被害人。（2011年第59题，2013年第61题）

4.使用假币罪与诈骗罪的关系。上述使用假币的方法中，有些同时触犯了诈骗罪，

例如，购物、消费；有些没有同时触犯诈骗罪，例如，冒充赃款上交、对机器使用。所以，使用假币罪与诈骗罪不是法条竞合关系，但可以想象竞合。

5.使用假币罪与盗窃罪的联系。

（1）盗窃普通财物，实际上窃得假币，又持有、使用的，定盗窃罪和持有、使用假币罪，并罚。

% Source PDF physical page 341; see audit/source_coverage.json.
（2）通过自动存取款机将假币存入银行，又从自动取款机取出真币的，构成使用假币罪和盗窃罪，因为是两个行为触犯两个法益，应并罚。（2009年第61题）

[总结］货币犯罪内部的罪数关系

\InsertDiagram{figures/manual_table_p341_1.tex}

\subsection*{二、洗钱罪}

第191条为掩饰、隐瞒\CJKunderwave{毒}品犯罪、\CJKunderwave{黑}社会性质的组织犯罪、\CJKunderwave{恐}怖活动犯罪、\CJKunderwave{走}私犯罪、\CJKunderwave{贪}污贿赂犯罪、破坏\CJKunderwave{金}融管理秩序犯罪、\CJKunderwave{金}融诈骗犯罪的所得及其产生的收益的来源和性质，有下列行为之一的，没收实施以上犯罪的所得及其产生的收益，处五年以下有期徒刑或者拘役，并处或者单处罚金；情节严重的，处五年以上十年以下有期徒刑，并处罚金：

\subsubsection*{（一）提供资金帐户的；}

\subsubsection*{（二）将财产转换为现金、金融票据、有价证券的；}

\subsubsection*{（三）通过转帐或者其他支付结算方式转移资金的；}

\subsubsection*{（四）跨境转移资产的；}

\subsubsection*{（五）以其他方法掩饰、隐瞒犯罪所得及其收益的来源和性质的。}

单位犯前款罪的，对单位判处罚金，并对其直接负责的主管人员和其他直接责任人员，依照前款的规定处罚。［《刑法修正案（十一）》修订]

1.七种上游犯罪的具体范围（毒-黑-恐-走-贪-金-金）

（1）毒品犯罪，是指刑法分则第六章第七节所规定的所有犯罪。

（2）黑社会性质的组织犯罪。这不仅包括第294条组织、领导、参加黑社会性质组织罪这个具体罪名，还包括以黑社会性质组织为主体实施的各种犯罪，例如，绑架罪、抢劫罪等。因此，虽然刑法未明文规定侵犯财产罪是洗钱罪的上游犯罪，但是，黑社会性质组织实施的侵犯财产罪，依然是洗钱罪的上游犯罪。（2011年第12题，2022年试题）

（3）恐怖活动犯罪，原理与黑社会性质的组织犯罪相同。

（4）走私犯罪，是指刑法分则第三章第二节所规定的全部走私犯罪。

（5）贪污贿赂犯罪，包括挪用公款罪（2007年第97题），不包括职务侵占罪。

% Source PDF physical page 342; see audit/source_coverage.json.
（6）破坏金融管理秩序犯罪，是指刑法分则第三章第四节所规定的犯罪。高利转贷罪是该节罪名（2009年第54题）。非法吸收公众存款罪是该节罪名。（2019年试题）

（7）金融诈骗犯罪，是指刑法分则第三章第五节所规定的犯罪。

[注意1]三个纯正的自然人犯罪：贷款诈骗罪、信用卡诈骗罪、有价证券诈骗罪。如果单位实施该三罪，按照单位的领导个人构成该三罪论处。其犯罪所得属于洗钱罪的对象。

[注意2］不包括合同诈骗罪，因为该罪属于扰乱市场秩序罪，而不属于金融诈骗罪。

2.上游犯罪的认定，只要有证据证明存在上游犯罪事实即可。在此前提下，有下列情形，不影响洗钱罪的认定：（1）上游犯罪尚未依法裁判；（2）上游犯罪行为人逃匿、死亡；（3）因罪数关系，以其他罪名论处；（4）上游犯罪超过追诉时效，洗钱罪没有超过追诉时效。（2011年第12题，2022年试题）

3.上游犯罪的认识错误

行为人只要认识到是七种犯罪所得及其收益，就具有洗钱罪的故意。即使行为人在这七种犯罪所得及其收益之间产生对象认识错误，也不影响洗钱罪故意的成立。例如，行为人将上游的毒品犯罪所得误认为是贪污犯罪所得而实施洗钱行为的，不影响洗钱罪的成立。

4.本罪与上游犯罪的帮助犯的区分：事先有无通谋。有，则属于上游犯罪的帮助犯；没有，则是洗钱罪。例1，甲为了贩毒，将购买毒品的款项，交给银行人员乙，让乙汇到境外，用于进货。乙照办。乙是贩卖毒品罪的帮助犯。例2，甲将贩毒赚的钱交给银行人员乙，让乙汇到境外。乙照办。乙构成洗钱罪。

5.罪数问题。（1）本罪与掩饰、隐瞒犯罪所得罪是法条竞合关系，按照特殊法优于一般法规则，以本罪论处。（2）本罪与窝藏、转移、隐瞒毒赃罪（第349条）想象竞合，择一重罪论处。\footnote{2024年8月20日最高人民法院、最高人民检察院《关于办理洗钱刑事案件适用法律若干问题的解释》。}

6.对“自洗钱行为”按照洗钱罪处罚。在《刑法修正案（十一）》之前，自己实施上游犯罪，然后自已洗钱，根据期待可能性原理，属于不可罚的事后行为。按照《刑法修正案（十一）》的规定，对这种自洗钱行为也应按照洗钱罪处罚。

（1）对上游犯罪和自洗钱构成的洗钱罪，应并罚。

（2）洗钱罪与上游犯罪能否想象竞合？少数说认为，上游犯罪与洗钱罪可以同时触犯，想象竞合。多数说认为，二者是上下游关系，不能想象竞合。

例1，甲欲向乙行贿，乙让甲直接将贿赂款汇到乙的境外账户。少数说认为，甲构成行贿罪与洗钱罪的想象竞合，乙构成受贿罪和洗钱罪的想象竞合，择一重罪论处。多数说认为，甲仅构成行贿罪，因为只有在上游犯罪既遂之后才可能构成洗钱罪。（2021年试题）

例2，官员甲为了实施贪污，将公款直接汇往境外账户。少数说认为，甲构成贪污罪和洗钱罪的想象竞合，择一重罪论处。多数说认为，甲仅构成贪污罪。

例3，甲要贩毒，让乙提供支付宝账户，用于接收货款（犯罪所得），乙照办。乙的账户收到后，乙转给甲。少数说认为，乙的一个提供结算账户的行为，一方面构成贩卖毒品罪的帮助犯，另一方面也构成洗钱罪，属于为自己和共犯甲洗钱，想象竞合，择一重罪论处。多数说认为，乙的行为仅构成贩卖毒品罪的帮助犯。

% Source PDF physical page 343; see audit/source_coverage.json.
（3）掩饰、隐瞒犯罪所得罪依然遵守期待可能性原理，也即，行为人实施上游财产犯罪后，实施掩饰、隐瞒行为（如销赃），属于不可罚的事后行为。

例如，甲贩毒赚了50万元现金，甲将其藏在地下室。甲的掩饰手段没有破坏金融管理秩序，不构成洗钱罪，属于普通的掩饰、隐瞒犯罪所得，属于不可罚的事后行为。

7.处罚：“没收上游犯罪的所得及其收益”。这里“没收”包括两个内容：一是，上游犯罪没有被害人的（例如，受贿罪、贩卖毒品罪），应当追缴上游犯罪所得及其收益。二是，上游犯罪有被害人的（例如，贪污罪、金融诈骗罪），应当将上游犯罪所得返还被害人，犯罪所得的收益上缴国库。

\subsection*{三、贷款类犯罪}

\subsubsection*{(一）高利转贷罪（第175 条)}

这是指，以转贷牟利为目的，套取金融机构信贷资金高利转贷他人，违法所得数额较大的行为。

1.本罪与贷款诈骗罪的区分。（2008年第11题）

（1）本罪以转贷牟利为目的，仍有向银行还本付息的意思。

（2）贷款诈骗罪（第193条）具有非法占有的目的，没有向银行还本付息的意思。

2.贷款时没有转贷牟利目的，正常贷款后产生转贷牟利目的并转贷的，不构成本罪。

3.变相高利转贷。

（1）行为人开始就有转贷牟利目的，套取金融机构信贷资金后，表面上将该部分资金用于生产经营，但又将自有资金高利借贷他人，违法所得数额较大的，应认定为本罪。

（2）行为人套取金融机构的信贷资金，高利借贷给名义上有合资合作关系但实际上并不参与经营的企业，违法所得数额较大的，也应认定为本罪。

\subsubsection*{（二）骗取贷款罪（第175条之一）}

\InsertDiagram{figures/manual_table_p343_1.tex}

\subsubsection*{（三）违法发放贷款罪（第186 条）}

本罪是指，银行或者其他金融机构的工作人员或者单位，违反国家规定发放贷款，数额巨大或者造成重大损失的行为。（2025年试题）

% Source PDF physical page 344; see audit/source_coverage.json.
1.这里的金融机构包括农村信用社，但不包括小额贷款公司。

2.本罪是故意犯罪。（1）如果贷款人欺骗银行人员，则银行人员不构成本罪。贷款人构成骗取贷款罪或贷款诈骗罪。（2）贷款人没有欺骗银行人员，而是唆使银行人员违法发放贷款，则银行人员构成本罪，但是贷款人不定本罪的共犯，这是立法者有意缩小处罚范围。

\subsection*{四、伪造、变造金融票证罪（第177条）}

本罪与四个金融诈骗犯罪是上下游的牵连关系，择一重罪论处。

\InsertDiagram{figures/manual_p344_1.tex}

例如，乙单位欠甲10万元，甲反复催讨后，乙单位开出10万元（没填大写）的现金支票给甲。甲持支票向银行询问乙单位账户上的存款数额，银行回答有31万元。于是甲回家将自己的9万元存入乙的账户，然后将支票数额改为40万元，从银行取出40万元。对甲的行为应认定为变造金融票证罪和票据诈骗罪。

\subsection*{五、妨害信用卡管理罪和窃取、收买、非法提供信用卡信息罪}

第177条之一有下列情形之一，妨害信用卡管理的，处三年以下有期徒刑或者拘役，并处或者单处一万元以上十万元以下罚金；····

\subsubsection*{（一）明知是伪造的信用卡而持有、运输的，或者明知是伪造的空白信用卡而持有、}

运输，数量较大的；

\subsubsection*{（二）非法持有他人信用卡，数量较大的；}

\subsubsection*{（三）\CJKunderwave{使用虚假的身份证明骗领信用卡}的；}

\subsubsection*{（四）出售、购买、为他人提供伪造的信用卡或者以虚假的身份证明骗领的信用卡的。}

\CJKunderwave{窃取、收买或者非法提供他人信用卡信息资料}的，依照前款规定处罚。

银行或者其他金融机构的工作人员利用职务上的便利，犯第二款罪的，从重处罚。

\section{金融诈骗罪}

\subsection*{一、信用卡诈骗罪}

信用卡诈骗罪是五星级罪名，涉及内容较多，需要专门讲解。复习方法是，先展开知识点，理解模型，最后记住三句话。

第196条\CJKunderwave{有下列情形之一，进行信用卡诈骗活动}，数额较大的，处五年以下有期徒刑或者拘役，并处二万元以上二十万元以下罚金；··

\subsubsection*{（一）\CJKunderwave{使用伪造的信用卡，或者使用以虚假的身份证明骗领的信用卡的}；}

\subsubsection*{（二）使用作废的信用卡的；}

% Source PDF physical page 345; see audit/source_coverage.json.
\subsubsection*{（三）\CJKunderwave{冒用他人信用卡的}；}

\subsubsection*{（四）\CJKunderwave{恶意透支的}。}

前款所称恶意透支，\CJKunderwave{是指持卡人以非法占有为目的，超过规定限额或者规定期限透}\CJKunderwave{支，并且经发卡银行催收后仍不归还的行为}。

\CJKunderwave{盗窃信用卡并使用的}，依照本法第二百六十四条（盗窃罪）的规定定罪处罚。

前提知识

\subsubsection*{(一)}

1.行为对象。“信用卡”。根据立法解释的规定，本罪的“信用卡”既包括具有透支功能的信用卡，也包括无透支功能的储蓄卡、借记卡。所以本罪的“信用卡”就是指银行卡。若是储蓄卡，则被害人是卡主。若是可透支的信用卡，则直接被害人是发卡银行。

2.行为方式。非法使用信用卡。按照原理，对人使用，欺骗人，构成信用卡诈骗罪。对机器使用，不存在欺骗人，则构成盗窃罪。但是，根据司法解释的特殊规定，非法使用信用卡，无论是对人使用，还是对机器使用，一律定信用卡诈骗罪。\footnote{2018年12月1日最高人民法院、最高人民检察院《关于办理妨害信用卡管理刑事案件具体应用法律若干问题的解释》第5条规定：“冒用他人信用卡”，包括以下情形：（1）拾得他人信用卡并使用的；（2）骗取他人信用卡并使用的；（3）窃取、收买、骗取或者以其他非法方式获取他人信用卡信息资料，并通过互联网、通讯终端等使用的；

（4）其他冒用他人信用卡的情形。2008年5月7日最高人民检察院《关于拾得他人信用卡并在自动柜员机（ATM机）上使用的行为如何定性问题的批复》规定：\CJKunderwave{拾得}他人信用卡并在ATM机上使用的行为，以信用卡诈骗罪论处。}这是一刀切的规定。做题时应按照司法解释作答。

\subsubsection*{（二）行为类型一：非法取得+使用假卡}

1.伪造信用卡（定伪造金融票证罪）+使用（定信用卡诈骗罪）=牵连犯，择一重罪论处，定信用卡诈骗罪。

例如，甲在ATM机上安装了读卡器，读取到卡主乙的银行卡信息，然后复制一张卡，这种复制卡的行为属于伪造卡，然后在ATM机上取款或对人使用。甲先构成窃取信用卡信息罪，然后构成伪造金融票证罪，再构成信用卡诈骗罪，三罪属于牵连犯关系，最终定信用卡诈骗罪。

2.以虚假身份证明骗领信用卡（定妨害信用卡管理罪）+使用（定信用卡诈骗罪）：牵连犯，择一重罪论处，定信用卡诈骗罪。

例如，甲用假身份证从银行骗领了一张可透支消费的信用卡，然后在商场购物，透支消费。甲先构成妨害信用卡管理罪，后构成信用卡诈骗罪，两罪是牵连犯关系，最终定信用卡诈骗罪。（2017年第58题，2019年试题）

\subsubsection*{（三）行为类型二：非法取得+冒用真卡}

1.非法取得行为的地位

在前文“财产犯罪的数额”中讲解过，存有3000元的银行卡属于价值数额不大、但值得刑法有限保护的财物。普通盗窃、诈骗、侵占、敲诈、抢夺银行卡，不使用，不构成犯罪。但是，抢劫、特殊类型盗窃（入户盗窃、多次盗窃、携带凶器盗窃、扒窃）银行卡，即使不使用，也构成抢劫罪、盗窃罪。既然一般的非法取得行为本身不构成犯罪，那么只有非法使用（冒用）了银行卡，处分了卡里资金才构成犯罪。

2.非法取得行为的种类一：诈骗、侵占、抢夺、敲诈勒索。

% Source PDF physical page 346; see audit/source_coverage.json.
[结论］\CJKunderwave{侵占、诈骗、抢夺、敲诈信用卡+使用=信用卡诈骗罪}。

例1，甲捡拾了乙遗失的银行卡，背面有密码。甲在ATM机上取款。甲构成信用卡诈骗罪。（2015年第57题）

例2，乙出国，让甲保管银行卡，背面有密码。甲擅自在商场刷卡消费。甲构成信用卡诈骗罪。（2013年第15题）

3.非法取得行为的种类二：抢劫。

[结论］\CJKunderwave{抢劫+使用=抢劫罪}。抢劫信用卡，即使不使用，也应定抢劫罪，按照情节量刑；如果使用，根据司法解释，仍然一律定抢劫罪，使用的数额就是抢劫的数额。\footnote{2016年1月6日最高人民法院《关于审理抢劫刑事案件适用法律若干问题的指导意见》。}（2025年试题）

例如，甲抢劫到乙的银行卡，逼迫乙说出密码，然后在ATM机上取走1万元。甲构成抢劫罪，数额按照1万元计算。（2006年第53题）

4.非法取得行为的种类三：普通盗窃。

[结论]\CJKunderwave{盗窃+（对人或机器）使用=盗窃罪}。刑法第196条第3款规定：\CJKunderwave{盗窃信用卡}\CJKunderwave{并使用的，以盗窃罪论处}。首先，普通盗窃信用卡的行为本身不构成盗窃罪，只有使用才构成犯罪。其次，按照原理，对机器使用，原本就定盗窃罪。对此，第196条第3款的结论是定盗窃罪。就此而言，该款规定是注意规定。按照原理，对人使用，骗人，原本应定信用卡诈骗罪，但是第196条第3款的结论是定盗窃罪，这意味着将信用卡诈骗罪拟制为盗窃罪，属于法律拟制。

例如（2007年第57题），下列说法是否正确：“丁盗窃信用卡并使用的行为，属于盗窃罪与信用卡诈骗罪的吸收犯”。官方答案：该说法错误。这是因为，该说法认为，丁盗窃信用卡的行为构成盗窃罪，非法使用行为构成信用卡诈骗罪，然后前者吸收了后者，因此最终定盗窃罪。然而，丁盗窃卡的行为本身不构成盗窃罪。最终定盗窃罪的依据不是盗窃行为，而是使用行为。法条在此将使用行为也即信用卡诈骗罪拟制为盗窃罪。

[巩固练习]题1（2011年第58题），下列说法是否错误：“即使没有刑法第196条第3款的规定，对于盗窃信用卡并在ATM机取款的行为，也能认定为盗窃罪”。官方答案：该说法正确。这是因为，盗窃卡的行为不构成盗窃罪，在机器上使用卡的行为本身就构成盗窃罪。此时，第196条第3款属于注意规定，也即删掉该款规定，在机器上使用卡的行为也能定盗窃罪。

题2（2017年第58题），下列说法是否错误：“刑法规定，盗窃信用卡并使用的，以盗窃罪论处。与此相应，拾得信用卡并使用的，就应以侵占罪论处”。官方答案：该说法错误。该说法认为，盗窃信用卡并使用定盗窃罪，依据的是盗窃卡的行为，以此类推，侵占信用卡并使用就应定侵占罪。然而，盗窃信用卡并使用定盗窃罪，依据的是使用行为，而非盗窃卡的行为。

[总结]命题套路

（1）\CJKunderwave{套路一}：判断非法取得行为的性质。

例如，甲外出几天，让朋友乙打扫房间，乙打扫时捡到甲的一张银行卡，乙事前知道密码，便私下去ATM机取出1万元。由于甲还在占有自己的卡，乙的捡拾不属于侵占遗

% Source PDF physical page 347; see audit/source_coverage.json.
忘物，属于盗窃卡的行为，所以应适用“盗窃信用卡并使用的，以盗窃罪论处”。

（2）\CJKunderwave{套路二}：共犯平台与认识错误。（2023年试题）

例如（2019年试题），甲盗窃了丙的一张信用卡，告诉乙是捡来的，让乙在商场购物。乙在商场冒用该信用卡购买了3万元商品。第一，乙的这种认识错误并不重要，因为乙至少认识到这是他人的卡，自已在非法使用，对人冒用，构成信用卡诈骗罪的实行犯。第二，甲首先可构成信用卡诈骗罪的教唆犯（指使）和帮助犯（提供卡），但是根据刑法规定，“盗窃信用卡并使用，即使对人使用，本应定信用卡诈骗罪，也拟制为盗窃罪”，因此最终对甲定盗窃罪。“盗窃+使用，定盗窃罪”，既包括亲自偷、亲自用，也包括亲自偷、让别人用。

\subsubsection*{（四）行为类型三：恶意透支}

恶意透支，是指持卡人以非法占有为目的，超过规定限额或者规定期限透支，并且经发卡银行两次有效催收后超过3个月仍不归还。\footnote{2018年12月1日最高人民法院、最高人民检察院《关于办理妨害信用卡管理刑事案件具体应用法律若干问题的解释》。}

1.卡的范围仅限于具有透支功能的信用卡，不包括没有透支功能的储蓄卡（借记卡）。

2.“恶意”透支，是指在透支时就有非法占有目的。如果透支时没有非法占有目的（具有归还意思），就不属于恶意透支。（2017年第58题）

3.“经发卡银行两次有效催收后超过3个月仍不归还”是客观处罚条件。实施完恶意透支，便触犯信用卡诈骗罪，但是此时立法者为了缩小打击面，给行为人一个机会“赶紧还款”，如果“经发卡银行两次有效催收后超过3个月仍不归还”，则启动刑罚处罚。该规定的性质类似于逃税罪第4款规定的排除处罚事由。

4.行为主体是卡主（合法持卡人）。这是指以自己的真实身份证申领到信用卡的持卡人。只有卡主，才能享受“催收程序”这个机会待遇。

例1，甲捡拾到乙的信用卡，恶意透支，定信用卡诈骗罪，属于“冒用他人信用卡”，不属于此处“恶意透支型信用卡诈骗罪”，不能享受“催收程序”这个机会待遇，因为银行不可能向甲催收。

例2，甲盗窃到乙的信用卡，恶意透支，属于“盗窃+使用”，定盗窃罪。

\subsubsection*{（五）银行卡信息资料}

[结论］以非法方式获取他人银行卡信息资料+（无论对人、机器、网络）使用=信用卡诈骗罪。这是2018年及2009年司法解释的一刀切规定。\footnote{2018年12月1日最高人民法院、最高人民检察院《关于办理妨害信用卡管理刑事案件具体应用法律若干问题的解释》。2009年12月16日最高人民法院、最高人民检察院《关于办理妨害信用卡管理刑事案件具体应用法律若干问题的解释》。}考试应当按照司法解释来答题。

例如，甲盗窃了乙的信用卡信息资料，通过网上银行（无人工参与，均是电脑系统处理），将乙账户里的资金划到自已账户。按照正常原理，这属于盗划行为，构成盗窃罪。但是按照司法解释，应定信用卡诈骗罪。

[巩固练习]题1（2016年第17题），下列说法是否正确：“丁侵入银行计算机信息系统，将刘某存折中的5万元存款转入自已的账户。对丁应以诈骗罪论处”。官方答案：该说

% Source PDF physical page 348; see audit/source_coverage.json.
法错误。该题故意绕开上述司法解释，针对的是存折，不是银行卡。因此，不需要遵守该司法解释，应按照正常原理处理。由于不存在对人使用、欺骗人的问题，因此构成盗窃罪。

题2（2019年试题），乙在工商银行领取银行卡及相配套的U盾，\footnote{U盾是工行推出的用于网上银行电子签名和数字认证的工具。使用U盾在登录个人网上银行之后，如需办理转账、汇款等支付业务，输入U盾密码，并经银行系统验证无误，即可完成支付业务。}银行大厅经理甲（非国家工作人员）在假意指导乙如何使用U盾时偷换了乙的U盾，并骗乙：“只能在一周后使用U盾。”乙信以为真。次日，甲使用乙的U盾将乙网上银行账户内的3万元转入自已的网上银行账户。甲盗窃了乙的U盾，就盗窃了乙的银行卡信息资料，通过网上银行，将乙账户里的资金划到自已账户。按照司法解释，对甲定信用卡诈骗罪。需要说明的是，甲窃取他人信用卡信息资料，构成窃取信用卡信息罪（第177条之一），与此后的信用卡诈骗罪是手段与目的的牵连犯关系，择一重罪论处，定信用卡诈骗罪。

\subsubsection*{（六）微信、支付宝的信息资料}

1.盗用他人微信、支付宝的信息资料

微信、支付宝账户不等于银行卡账户，微信、支付宝的信息资料不等于银行卡信息资料。因此，不需要适用2018年司法解释的特殊规定（一律定信用卡诈骗罪）。

[结论】盗用微信、支付宝的信息资料，取得主人账户内资金，一律定盗窃罪。\footnote{参见张明楷：《侵犯人身罪与侵犯财产罪》，北京大学出版社2021年版，第221、232页。}

例1（微信钱包、支付宝账户内的余额），甲获取了乙手机上微信的支付密码。第一种情形，盗划，也即将乙微信钱包的1万元转入自已的银行账户。第二种情形，盗刷，也即消费掉乙微信钱包的1万元。甲的盗划行为、盗刷行为，一律定盗窃罪。\footnote{可能有同学问：若对人（收银员）使用，为何不定诈骗罪？理由是，使用银行卡时，一般需要使用人签字，表明收银员对使用者的身份有审核义务。而使用微信、支付宝，都是扫二维码，一般不需要使用人签字，表明收银员对使用者的身份没有审核义务。没有审核义务，就不存在受骗问题。}（2020年试题）

例2（微信、支付宝绑定的银行账户），甲获取了乙手机上微信的支付密码，发现乙的微信钱包中没钱，但是乙的微信绑定了乙的一个银行卡账户，该账户里有1万元。第一种情形，盗划，也即甲用微信支付密码将该银行账户内的1万元转到自己的银行账户。第二种情形，盗刷，也即甲用微信支付密码消费掉这1万元。关于甲的盗划行为、盗刷行为：

（1）多数说认为，定盗窃罪。理由：虽然1万元是银行账户的资金，但是甲非法使用的是微信支付密码，微信支付密码不等于银行卡信息资料，因此在此不需要适用2018年司法解释的特殊规定（一律定信用卡诈骗罪）。

（2）少数说认为，定信用卡诈骗罪。理由：既然使用微信支付密码能够取出银行卡里的资金，那么微信支付密码便等同于银行卡支付密码，因此适用2018年司法解释的特殊规定，定信用卡诈骗罪。（2022年主观题、2023年客观题）

2.蚂蚁花呗

这是蚂蚁金服公司推出的一款透支消费的借贷服务，申请开通后，可获得数万元的消费透支额度，先消费，后还款。刑法界主流观点认为，虽然花呗的这种透支消费服务类似

% Source PDF physical page 349; see audit/source_coverage.json.
于可以透支的信用卡服务，但是花呗在法律性质上不同于信用卡。\footnote{主要理由是，根据刑法第196条第2款规定，可以透支的信用卡是由发卡银行发行的，而蚂蚁金服公司不属于银行。参见张明楷：《合同诈骗罪行为类型的边缘问题》，载《东方法学》2020年第1期。}因此，花呗的信息资料不属于信用卡信息资料，对此不需要适用2018年司法解释（一律定信用卡诈骗罪）。

消费信贷也是贷款行为。蚂蚁金服公司虽然不是银行，但属于“其他金融机构”。贷款诈骗罪保护其他金融机构。因此，盗用他人名义，向蚂蚁金服公司申请消费类贷款，构成贷款诈骗罪。

[结论］冒用他人蚂蚁花呗，取得蚂蚁公司的垫付资金，一律定贷款诈骗罪。

例如（2019年试题），甲获取了乙手机上的支付宝密码，使用里面的“花呗”在商场购买了3万元商品。甲冒用“花呗”，不构成信用卡诈骗罪，而构成贷款诈骗罪。

3.蚂蚁借呗

这是蚂蚁金服公司推出的一款借贷服务，行为人申请贷款，蚂蚁借呗批准后，将款项打人行为人的指定账户，行为人日后还款。借呗与花呗均属于向蚂蚁金服公司借款，然后还款。二者区别在于，借呗将出借款项打入行为人的账户；花呗将出借款项（垫付款）打人商家的账户。借呗的借贷服务中，不存在信用卡的问题，因此不需要适用2018年司法解释，按照正常原理处理。

[结论］冒用他人蚂蚁借呗，取得蚂蚁公司的放贷资金，一律定贷款诈骗罪。

例如，甲获取了乙手机上的支付宝密码，使用里面的“借呗”申请贷款2万元，蚂蚁公司审核批准后，将款项打入甲的指定账户。甲冒用“借呗”，构成贷款诈骗罪。

引申情形：如果甲让蚂蚁公司将款项打人乙的支付宝账户，那么，甲仍构成贷款诈骗罪（多数说）。此后，第一种情形，甲悄悄从乙的账户划走该资金，甲构成盗窃罪。盗窃罪与贷款诈骗罪不需并罚，重罪吸收轻罪。第二种情形，甲欺骗乙，让乙转给了自己，甲构成诈骗罪。诈骗罪与贷款诈骗罪不需并罚，重罪吸收轻罪。

[总结1]\CJKunderwave{实体卡}

（1）使用假卡

\InsertDiagram{figures/manual_table_p349_1.tex}

(2）冒用真卡

\InsertDiagram{figures/manual_table_p349_2.tex}

% Source PDF physical page 350; see audit/source_coverage.json.
[总结2］\CJKunderwave{银行卡信息资料}

\InsertDiagram{figures/manual_table_p350_1.tex}

[提示］可以再总结，凡是非法使用银行卡（包括实体卡和信息资料），一般都定信用卡诈骗罪，除了两个情形：抢劫+使用，定抢劫罪：盗窃+使用，定盗窃罪。

[总结3］微信、支付宝的信息资料

\InsertDiagram{figures/manual_table_p350_2.tex}

[记忆】分析一大堆，最后记住三句话：

\InsertDiagram{figures/manual_p350_3.tex}

\subsection*{二、贷款诈骗罪（第193条）}

该罪是指，以非法占有为目的，诈骗银行或者其他金融机构的贷款，数额较大的行为。

1.行为主体：只能是自然人。单位实施贷款诈骗的，追究领导个人的贷款诈骗罪。

2.罪名区分。（1）贷款诈骗罪与骗取贷款罪的区分：前者有非法占有目的，后者没有非法占有目的。（2）贷款诈骗罪与高利转贷罪的区分：前者具有非法占有目的；后者没有非法占有的目的，但有转贷牟利的目的。

3.虚假担保问题（两头骗）。

模型1，甲想实施贷款诈骗，欺骗乙，让乙将乙的汽车作为质押物，带着非法占有贷款的目的，向银行贷到款，然后跑路。贷款到期，银行未能收到本息，便向乙行使质押权，获得了追偿，挽回了损失。甲对乙构成诈骗罪，诈骗的对象是财产性利益（担保物权）。甲对银行是否构成贷款诈骗罪，主要看银行是不是被害人，有无财产损失的危险。

（1）多数说：银行有财产损失的危险，是被害人，甲对银行构成贷款诈骗罪。理由：第一，甲有非法占有目的，银行就有财产损失的危险。第二，根据个别财产说，银行放贷的目的是到期收回本息。因此，财产损失的认定时间应以贷款到期时为准。此时银行未能收到本息，就有实际财产损失。至于事后银行追偿，那是挽回损失。挽回损失以遭受财产损失为前提。第三，银行能否顺利行使质押权，存在风险，因此银行至少存在损失风险。在罪数上，诈骗乙是手段行为，诈骗银行是目的行为，二者是牵连犯关系，择一重罪论处。

（2）少数说：银行没有财产损失的危险，不是被害人，甲对银行不构成贷款诈骗罪。理由：第一，虽然甲有非法占有目的，但是甲提供了足额担保，保障清偿，所以银行没有财产损失的危险。第二，根据整体财产说，银行放贷的目的是营利，财产损失的认定时间

% Source PDF physical page 351; see audit/source_coverage.json.
应以追偿后为准，看有无实际财产损失。既然银行获得了追偿，便表明没有实际财产损失。（2022年主观题）

模型2，甲诈骗到乙的汽车，然后谎称是自己的汽车，作为质押物提供给银行，带着非法占有贷款的目的，向银行贷到款，然后跑路。贷款到期，银行未能收到本息，行使质押权，获得追偿。甲对乙构成诈骗罪。多数说认为，甲对银行构成贷款诈骗罪。两罪并罚。（2019年主观题）

4.共犯问题。

模型1，普通公民甲与银行的丙（部门审核人员，国家工作人员）勾结，以非法占有为目的，共同欺骗分管领导乙（贷款最终决定者），使其产生认识错误并核准贷款。甲同时触犯贷款诈骗罪的实行犯和贪污罪的帮助犯，\footnote{如果加个条件“丙在贷款诈骗中起到支配作用”，则丙在帮助犯的基础上就升级为共同正犯。}想象竞合，择一重罪论处。丙同时触犯贪污罪的实行犯和甲的贷款诈骗罪的帮助犯，想象竞合，择一重罪论处。

模型2，普通公民甲与银行的乙（领导，贷款最终决定者，国家工作人员）勾结，欺骗下级部门审核人员丙，丙被骗后，将甲提供的虚假材料递交给乙，乙批准了贷款。虽然丙受骗，但丙不是处分人，因此甲不构成贷款诈骗罪。乙构成贪污罪的实行犯，甲构成贪污罪的帮助犯。

\subsection*{三、保险诈骗罪（第198条）}

保险诈骗罪，是指投保人、被保险人、受益人，以使自己或者第三人获取保险金为目的，采取虚构保险标的、制造保险事故等方法，骗取保险金的行为。

\InsertDiagram{figures/manual_table_p351_1.tex}

% Source PDF physical page 352; see audit/source_coverage.json.
\ContinuedTable{figures/manual_table_p352_1.tex}
\SourceFootnote{\textbf{表注①}\quad 对此不认定牵连犯的理由是，虽然具有手段与目的关系，但是，不是常见常发的情形。}

[总结]医保骗保问题\footnote{2024年2月28日最高人民法院、最高人民检察院、公安部《关于办理医保骗保刑事案件若干问题的指导意见》。}

1.行为人以非法占有为目的，骗取医疗保障基金支出的，构成诈骗罪。

2.定点医药机构（医疗机构、药品经营单位）以非法占有为目的，骗取医疗保障基金支出的，对直接责任人员定诈骗罪。

3.盗窃他人医疗保障凭证（社会保障卡等），并盗刷个人医保账户资金，构成盗窃罪。

\subsection*{四、集资类犯罪}

\InsertDiagram{figures/manual_table_p352_2.tex}

% Source PDF physical page 353; see audit/source_coverage.json.
\ContinuedTable{figures/manual_table_p353_1.tex}

\section{危害税收征管罪}

[问题］狗蛋一直逃税，税务人员上门征税，狗蛋暴力抗税。后狗蛋还是补缴应纳税款，缴纳滞纳金，也接受了行政处罚。可不可以不追究狗蛋的逃税罪的刑事责任？（抗税案）

\subsection*{一、逃税罪}

第201条纳税人采取欺骗、隐瞒手段进行虚假纳税申报或者不申报，逃避缴纳税款数额较大并且占应纳税额百分之十以上的，处三年以下有期徒刑或者拘役，并处罚金；…

扣缴义务人采取前款所列手段，不缴或者少缴已扣、已收税款，数额较大的，依照前款的规定处罚。

对多次实施前两款行为，未经处理的，按照累计数额计算。

有\CJKunderwave{第一款}行为，经税务机关依法下达追缴通知后，\CJKunderwave{补缴应纳税款，缴纳滞纳金}，\CJKunderwave{已受}\CJKunderwave{行政处罚}的，\CJKunderwave{不予追究}刑事责任；\CJKunderwave{但是，五年内}因\CJKunderwave{逃避缴纳税款受过刑事处罚}或者\CJKunderwave{被税务}\CJKunderwave{机关给予二次以上行政处罚}的除外。

1.行为主体是特殊主体，即纳税人和扣缴义务人。

[提示］税务机关的工作人员与纳税人或者扣缴义务人相勾结，共同实施逃税行为的，成立逃税罪的共犯；行为同时触犯徇私舞弊不征、少征税款罪（第404条），想象竞合，择一重罪论处。

2.排除处罚事由（排除刑事责任的事由）。本条第4款是在构成犯罪的情况下排除刑事责任。（2009年第54题，2012年第61题）

（1）任何逃税案件，首先必须经过税务机关的处理。税务机关没有处理或者不处理的，司法机关不得直接追究行为人的刑事责任。

（2）如果税务机关只要求行为人补缴应纳税款，缴纳滞纳金，而没有给予行政处罚的，只要行为人补缴应纳税款和缴纳滞纳金，就不应追究其刑事责任。

（3）行为人只要处在行政法处理期间，就可享受该排除处罚事由的待遇。但一旦进入刑事司法程序，即使补缴应纳税款，缴纳滞纳金，接受行政处罚，也不能享受排除处罚事由的待遇。

（4）但书规定的“二次以上行政处罚”是指已经受到二次行政处罚，第三次再逃税的，才不能适用排除处罚事由的待遇。注意：因漏税而受行政处罚的，不包含在内。

（5）该排除处罚事由的待遇仅适用于纳税人，不适用于扣缴义务人的逃税行为，也不适用于抗税罪（第202条）、逃避追缴欠税罪（第203条）的行为人。例如，本节开头“抗税案”中，狗蛋不能享受这个待遇。

% Source PDF physical page 354; see audit/source_coverage.json.
\subsection*{二、骗取出口退税罪}

第204条以假报出口或者其他欺骗手段，骗取国家出口退税款，数额较大的，处五年以下有期徒刑或者拘役，并处骗取税款一倍以上五倍以下罚金；·.

\CJKunderwave{纳税人缴纳税款后，采取前款规定的欺骗方法，骗取所缴纳的税款的，依照本法第二}\CJKunderwave{百零一条（逃税罪）的规定定罪处罚；骗取税款超过所缴纳的税款部分，依照前款的规定}\CJKunderwave{处罚}。

[注意】本罪第2款对想象竞合进行数罪并罚。例如，甲先依法缴纳了税款100万元，然后采取骗取出口退税的方式骗回来120万元。一方面，甲的100万元税款等于没交，构成逃税罪。另一方面，甲对多余的20万元构成骗取出口退税罪。甲的一个骗回行为同时触犯两罪，想象竞合，本应择一重罪论处，但是依照该款规定，应数罪并罚。这是刑法典唯一一个对想象竞合进行数罪并罚的特殊规定。（2005年第10题，2008年第59题）

\subsection*{三、虚开发票罪与虚开增值税专用发票罪}

\subsubsection*{（一）虚开发票罪（第205条之一）}

1.该罪的发票是指普通发票。但是，增值税专用发票可以评价为普通发票。

2.该罪的保护法益是发票的公共信用。

3.该罪的成立条件是，带着非法使用目的，实施虚开行为。虚开行为是该罪唯一的实行行为。带着非法使用目的而虚开，才会对发票的公共信用造成危险。

\subsubsection*{（二）虚开增值税专用发票罪（第205条）\protect\footnotemark}\footnotetext{（1）增值税的缴纳流程。例如，第一，甲服装厂向乙布料厂购进一批布料，需向乙厂支付113万元，其中，货物价值是100万元，13万元是进项税额=100万\ensuremath{\times}13\%（增值税税率），是购买方（甲厂）向销售方（乙厂）支付的，销售方（出票方）开具增值税专用发票给购买方（受票方）。第二，甲厂将布料加工成衣服，出售给丙服装店，收到货款226万元，其中，货物价值是200万元，26万元是销项税额=200万\ensuremath{\times}13\%，是购买方（丙店）向销售方（甲厂）支付的，甲厂开具增值税专用发票给丙店。第三，甲厂应缴纳的增值税=销项税额（26万）-进项税额（13万）=13万。概言之，销项税额是甲厂（销售方）收取购买方（丙店）的税额，进项税额是甲厂（购买方）支付销售方（乙厂）的税额。

（2）虚开增值税专用发票，骗抵税款的流程。例如，第一，上述甲厂与丁公司合谋，虚构一个“甲厂购买丁公司100万元布料”的合同，丁公司不需要给甲厂发货，甲厂不需要给丁公司支付货款，而向丁公司支付一笔买票费用，丁公司给甲厂开具13万元进项税额的增值税专用发票。第二，甲厂向丙服装店真实销售衣服，收到货款226万元，其中26万是销项税额。第三，甲厂向税局称，应缴纳增值税=销项数额（26万）-进项税额（虚构的13万元）=13万。甲厂少缴纳13万元税款。甲厂属于让他人（丁公司）为自已虚开增值税发票。丁公司属于为他人虚开增值税发票。}

1.本罪的构成要件（2008年第59题，2024年试题）

虚开行为只是本罪的预备行为，带着骗抵税款的目的实施的骗抵行为，是本罪的实行行为。骗抵成功，造成国家税收损失，就构成既遂。未骗抵成功，则构成未遂。\footnote{少数说认为，带着骗抵税款的目的进行虚开，就成立本罪；虚开出来，就构成既遂。虚开行为是本罪的实行行为。极少数说认为，只要有虚开行为，就成立本罪；虚开出来，就构成既遂。虚开行为是本罪的实行行为。}

（1）带着骗抵税款的目的实施了虚开行为，未实施骗抵行为，一方面构成虚开增值税专用发票罪的犯罪预备（或预备阶段的犯罪中止），另一方面构成虚开发票罪，因为增值税专用发票可以评价为普通发票，该罪的实行行为是虚开行为；想象竞合，择一重罪论处，以虚开发票罪论处。

% Source PDF physical page 355; see audit/source_coverage.json.
（2）虚开时，没有带着骗抵税款的目的，虚开后带着骗抵税款目的实施骗抵行为，构成虚开增值税专用发票罪。\footnote{参见张明楷：《论虚开增值税专用发票罪的构造》，载《清华法学》2024年第4期；2024年3月20日最高人民法院、最高人民检察院《关于办理危害税收征管刑事案件适用法律若干问题的解释》。提示：本罪法条仅规定了虚开行为，没有规定骗抵税款行为。张明楷教授认为，骗抵税款行为是不成文的（隐藏的）构成要件要素。}

例1，甲、乙双方以相同的数额相互为对方虚开增值税专用发票，并且已按规定缴纳税款，没有骗抵税款的目的，不构成虚开增值税专用发票罪。

例2，甲为了虚增公司业绩，所虚开的增值税专用发票没有抵扣联的，不构成虚开增值税专用发票罪。

[巩固练习]题1（2024年试题），农民乙卖给甲公司300万元的蘑菇，但是乙没有办法给甲公司开具增值税专用发票。甲公司找到乙公司，请求帮忙开具300万元的增值税专用发票，乙公司同意并照办。由于甲公司与乙公司没有骗抵税款的目的，不构成虚开增值税专用发票罪。

题2，甲用乙公司名义对外签订销售合同，由乙公司收取货款、开具增值税专用发票。甲没有骗抵税款的目的，不构成虚开增值税专用发票罪。

2.认识错误。例如，甲以为自已虚开的是增值税专用发票，实际上虚开的是普通发票。对甲可认定为虚开发票罪。

3.根据刑法第210条，盗窃增值税专用发票，定盗窃罪；骗取增值税专用发票，定诈骗罪。

\section{侵犯知识产权罪}

\subsection*{一、侵犯著作权罪（第217条）}

该罪是指，\CJKunderwave{以营利为目的}，侵犯著作权或者与著作权有关的权利，违法所得数额较大或者有其他严重情节的行为。常考的行为类型：

1.未经著作权人许可，复制\CJKunderwave{发行}、通过信息\CJKunderwave{网络向公众传播}其作品。两种行为类型：一是“发行”，二是“网络传播”。

\InsertDiagram{figures/manual_table_p355_1.tex}

例1，甲带着营利目的，将盗版电影文件上传网上，供不特定人下载，属于发行，属于网络传播，构成本罪。

例2，甲带着营利目的，将盗版电影光盘卖给、租给不特定人，属于线下发行，构成本罪。

% Source PDF physical page 356; see audit/source_coverage.json.
例3，甲带着营利目的，将盗版电影文件上传网上，供不特定人在线观看，但不能下载，属于播放，属于网络传播，构成本罪。

例4，老板甲在KTV用VCD影碟机播放他人音乐作品给客人听，属于线下播放，不构成本罪。（2014年第51题）

2.未经著作权人或者与著作权有关的权利人许可，故意避开或者破坏权利人为其作品、录音录像制品等采取的保护著作权或者与著作权有关的权利的技术措施。例如，权利人对计算机软件做了加密保护措施，行为人擅自解密。

3.罪数。

（1）既实施侵犯行为，又实施销售行为，只定侵犯著作权罪。

（2）不实施侵犯行为，只销售他人的侵权作品，定销售侵权复制品罪。（2025年试题）

（3）既实施侵犯行为，又销售他人的侵权作品，定侵犯著作权罪、销售侵权复制品罪，并罚。

\subsection*{二、侵犯商业秘密罪}

第219条有下列侵犯商业秘密行为之一，情节严重的，处三年以下有期徒刑，并处或者单处罚金；情节特别严重的，处三年以上十年以下有期徒刑，并处罚金：

\subsubsection*{（一）以盗窃、\CJKunderwave{贿赂}、欺诈、胁迫、\CJKunderwave{电子侵入}或者其他不正当手段获取权利人的商业}

秘密的；

\subsubsection*{（二）披露、使用或者允许他人使用以前项手段获取的权利人的商业秘密的；}

\subsubsection*{（三）违反保密义务或者违反权利人有关保守商业秘密的要求，披露、使用或者允许}

他人使用其所掌握的商业秘密的。

明知前款所列行为，获取、披露、使用或者允许他人使用该商业秘密的，以侵犯商业秘密论。

本条所称权利人，是指商业秘密的所有人和经商业秘密所有人许可的商业秘密使用人。［《刑法修正案（十一）》修订]

1.《刑法修正案（十一）》修改了本罪的成立条件。原成立条件是“给权利人造成重大损失”。现成立条件为“情节严重”。也即，将本罪从结果犯修改为情节犯。

2.采取非法复制等方式获取商业秘密的，属于“盗窃”商业秘密。\footnote{2025年4月26日最高人民法院、最高人民检察院《关于办理侵犯知识产权刑事案件适用法律若干问题的解释》。}

3.《刑法修正案（十一）》增设“为境外窃取、刺探、收买、非法提供商业秘密罪”（第219条之一）。

\subsection*{三、假冒注册商标罪}

第213条未经注册商标所有人许可，\CJKunderwave{在同一种商品、服务上使用与其注册商标相同}\CJKunderwave{的商标}，情节严重的，处三年以下有期徒刑，并处或者单处罚金；·….…［《刑法修正案

\subsubsection*{（十一）》修订]}

1.《刑法修正案（十一）》增加了对服务商标的保护。由此，本罪的商标包括商品商标和服务商标。

% Source PDF physical page 357; see audit/source_coverage.json.
2.罪数。

（1）既实施假冒行为，又实施销售行为，只定假冒注册商标罪。

（2）不实施假冒行为，只销售他人的假冒商品，定销售假冒注册商标的商品罪。

（3）既实施假冒行为，又销售他人的假冒商品，定假冒注册商标罪、销售假冒注册商标的商品罪，并罚。

\section{扰乱市场秩序罪}

\subsection*{一、非法经营罪}

第225条违反国家规定，有下列非法经营行为之一，扰乱市场秩序，情节严重的，处五年以下有期徒刑或者拘役，并处或者单处违法所得一倍以上五倍以下罚金；··

\subsubsection*{（一）未经许可经营法律、行政法规规定的专营、专卖物品或者其他\CJKunderwave{限制买卖的物品}的；}

\subsubsection*{（二）\CJKunderwave{买卖进出口许可证}、进出口原产地证明以及其他法律、行政法规规定的\CJKunderwave{经营许}}

\CJKunderwave{可证}或者批准文件的；

\subsubsection*{（三）未经国家有关主管部门批准非法经营证券、期货、保险业务的，或者\CJKunderwave{非法从事}}

\CJKunderwave{资金支付结算业务}的；

\subsubsection*{（四）其他严重扰乱市场秩序的非法经营行为。}

本罪中的“非法经营”不是泛指一切非法经营活动，而是指该经营活动需要行政特别许可，但未经行政特别许可而非法经营。

\subsubsection*{（一）第一项行为要点}

1.烟草属于专营专卖物品。

2.食盐已经不是特许经营物品。非法经营食盐，不构成非法经营罪。将工业碘盐添加到食盐中，予以生产、销售，定生产、销售有毒、有害食品罪。生产、销售含碘量不符合标准的食盐，定生产、销售不符合安全标准的食品罪。（2020年试题）

3.非法倒卖陈化粮，不构成非法经营罪。

4.无证经营成品油，不构成非法经营罪。

5.个人收购、销售黄金，不构成非法经营罪。

\subsubsection*{（二）第三项行为的要点}

资金支付结算业务（简称支付业务），是指商业银行或者支付机构在收付款人之间提供的货币资金转移服务。从事该业务需要取得央行的许可证。\footnote{非法从事资金支付结算业务的情形有：（1）非法设立第三方支付平台。例如，非法设立类似支付宝的第三方支付平台，构成非法经营罪。（2）非法发行可跨地区、跨行业、跨法人使用的多用途预付卡。（3）使用受理终端或者网络支付接口等方法，以虚构交易、虚开价格、交易退款等非法方式向指定付款方支付货币资金。（4）非法为他人提供支票套现服务。所依据的司法解释有：2017年6月2日最高人民检察院公诉厅《关于办理涉互联网金融犯罪案件有关问题座谈会纪要》；2019年2月1日最高人民法院、最高人民检察院《关于办理非法从事资金支付结算业务、非法买卖外汇刑事案件适用法律若干问题的解释》。}（2009年第54题）

\subsubsection*{（三）第四项行为的扩张}

对第四项行为（兜底条款）应依据同类解释规则来解释。不过，许多司法解释扩张了第四项行为，由此扩张了非法经营罪的成立范围，归纳如下（以年份排列）：

% Source PDF physical page 358; see audit/source_coverage.json.
1.\CJKunderwave{非法买卖外汇}。

2.\CJKunderwave{经营非法出版物}。注意：非法经营罪与侵犯著作权罪的区别：犯罪对象不同。前者的出版物本身就没有合法的著作权，例如，非法传销的宣传册；后者的出版物是盗版制品，所盗版的对象是拥有著作权的合法出版物，例如，盗版的《百年孤独》，《百年孤独》是作者享有著作权的作品。（2009年第57题）

3.\CJKunderwave{擅自经营国际电信业务}。（2009年第57题）

4.\CJKunderwave{非法发行彩票}。（2014年第14题）

5.非法\CJKunderwave{使用POS机提现}。违反国家规定，使用销售点终端机具（POS机）等方法，以虚构交易、虚开价格、现金退货等方式向信用卡持卡人直接支付现金，以非法经营罪定罪处罚。\footnote{2018年12月1日最高人民法院、最高人民检察院《关于办理妨害信用卡管理刑事案件具体应用法律若干问题的解释》。}

6.非法从事\CJKunderwave{网络水军活动}。违反国家规定，以营利为目的，通过信息网络有偿提供删除信息服务，或者明知是虚假信息，通过信息网络有偿提供发布信息等服务，扰乱市场秩序，情节严重的，以非法经营罪论处。\footnote{2013年9月10日最高人民法院、最高人民检察院《关于办理利用信息网络实施诽谤等刑事案件适用法律若干问题的解释》。}

7.放高利贷。司法解释规定，\footnote{（1）上述中的“经常性地向社会不特定对象发放贷款”，是指2年内向不特定多人（包括单位和个人）以借款或其他名义出借资金10次以上。注意：仅向亲友、单位内部人员等特定对象出借资金，不得适用上述规定。（2）上述中的“情节严重”，是指以超过36\%的实际年利率实施符合上述非法放贷行为，具有下列情形之一：个人非法放贷数额累计在200万元以上的，单位非法放贷数额累计在1000万元以上的。参见2019年10月21日最高人民法院、最高人民检察院、公安部、司法部《关于办理非法放贷刑事案件若干问题的意见》。}\CJKunderwave{以营利为目的，经常性地向社会不特定对象发放贷}\CJKunderwave{款}，情节严重的，以非法经营罪论处。

8.擅自从事互联网上网服务，以非法经营罪论处。\footnote{2024年国务院《互联网上网服务营业场所管理条例》第27条。}

9.擅自设立出版物的出版、印刷、进口单位，擅自从事出版物的出版、印刷、进口、发行业务，以非法经营罪论处。\footnote{2024年国务院《出版管理条例》第61条。}

\subsection*{二、强迫交易罪}

第226条以暴力、威胁手段，实施下列行为之一，情节严重的，处三年以下有期徒刑或者拘役，并处或者单处罚金；·

\subsubsection*{（一）强买强卖\CJKunderwave{商品}的；}

\subsubsection*{（二）强迫他人提供或者接受\CJKunderwave{服务}的；}

\subsubsection*{（三）强迫他人参与或者退出\CJKunderwave{投标、拍卖}的；}

\subsubsection*{（四）强迫他人转让或者收购公司、企业的\CJKunderwave{股份、债券或者其他资产}的；}

\subsubsection*{（五）强迫他人参与或者退出\CJKunderwave{特定的经营活动}的。}

1.本罪与抢劫罪、敲诈勒索罪不是A与-A的对立排斥关系。一个行为符合强迫交易罪，并不意味着不构成抢劫罪、敲诈勒索罪，如果同时符合抢劫罪、敲诈勒索罪，属于想象竞合，择一重罪论处。例如，某导游将旅客带到某“友谊商店”，店主要求旅客必须高

% Source PDF physical page 359; see audit/source_coverage.json.
价购买某纪念商品，否则拿刀砍。店主同时触犯强迫交易罪和抢劫罪，想象竞合，择一重罪论处，定抢劫罪。

2.司法解释规定：以暴力、胁迫手段强迫他人借贷，属于“强迫他人提供或者接受服务”，情节严重的，以强迫交易罪追究刑事责任。以非法占有为目的，以借贷为名采用暴力、胁迫手段获取他人财物，以抢劫罪或者敲诈勒索罪追究刑事责任。（2014年第14题）

\subsection*{三、提供虚假证明文件罪（第229条）}

本罪是指，承担某些职责的中介组织的人员故意提供虚假证明文件，情节严重的行为。

1.行为主体。第一，不限于承担资产评估、验资、会计、审计、法律服务等职责的中介组织的人员，还包括其他行业领域的中介组织的人员。第二，不限于有资质的中介组织的人员，还包括无资质的中介组织的人员。例如，某企业向政府机关提交虚假的人才引进方面的证明文件，以便为他人落户。该企业构成本罪。（2023年试题）

2.本罪是故意犯罪。对应的过失犯罪叫“出具证明文件重大失实罪”。

3.罪数。犯本罪，又受贿的，择一重罪论处。受贿的罪名一般是非国家工作人员受贿罪。

\subsection*{四、合同诈骗罪}

1.普通诈骗罪（A）与合同诈骗罪（A+B）是一般法条与特殊法条的关系，A是指侵犯他人财产权，B是指扰乱了市场秩序。由于要求扰乱市场秩序，因此，这里的签订、履行合同应限于比较正式的经济合同。例如，货物买卖合同、借贷合同、运输合同、承包合同等。例如，甲得知乙（一般公民，并非从事经营活动的市场主体）有大量存款，欺骗乙：“我有一个很好的投资项目，你若投资10万元，3个月后返还20万元。”乙答应，二人签了投资协议。甲获得10万元后潜逃。由于甲的行为没有达到扰乱市场秩序的程度，对甲应以普通诈骗罪论处。

2.由于一些金融诈骗罪（如贷款诈骗罪、保险诈骗罪）需要签订合同，与合同诈骗罪是特殊法条与一般法条的法条竞合关系，优先适用特殊法条。例如（2017年第14题），利用合同诈骗保险金，应认定为保险诈骗罪。

\subsection*{五、组织、领导传销活动罪}

第224条之一组织、领导以推销商品、提供服务等经营活动为名，要求参加者以缴纳费用或者购买商品、服务等方式获得加入资格，并按照一定顺序组成层级，直接或者间接以发展人员的数量作为计酬或者返利依据，引诱、胁迫参加者继续发展他人参加，\CJKunderwave{骗取}\CJKunderwave{财物}，扰乱经济社会秩序的传销活动的，处五年以下有期徒刑或者拘役，并处罚金；情节严重的，处五年以上有期徒刑，并处罚金。

1.传销与合法销售的关键区别在于，传销是以发展人员的数量作为计酬或者返利依据。

2.条文中的“骗取财物”是指本罪的一种间接目的，并不要求有对应的诈骗行为，更不要求实际骗取到财物。
